\documentclass[11pt,a4paper]{article}

\usepackage{amsmath}
\usepackage{mathtools}
\usepackage{fontspec}
\usepackage{unicode-math}
\setmathfont{KpMath-Regular.otf}
\setmathfont{KpMath-Bold.otf}[version=bold]
\usepackage{microtype}
\usepackage[dvipsnames]{xcolor}
\usepackage[margin=24mm, top=28mm, bottom=28mm]{geometry}

\usepackage{fancyhdr}
\pagestyle{fancy}
\fancyhf{}
\fancyhead[L]{\small\itshape The Voxel \texorpdfstring{$T$}{T}-Matrix}
\fancyhead[R]{\small\itshape 19 September 2026}
\fancyfoot[C]{\thepage}
\renewcommand{\headrulewidth}{0.4pt}

\setcounter{secnumdepth}{3}
\usepackage{booktabs}
\usepackage{array}
\usepackage[font=small, labelfont=bf, labelsep=period]{caption}
\usepackage{setspace}
\setstretch{1.12}
\setlength{\parskip}{0.4em}
\usepackage[colorlinks=true, linkcolor=blue!60!black, urlcolor=blue!60!black]{hyperref}

\newcommand{\bI}{\mathbf{I}}
\newcommand{\bP}{\mathbf{P}}
\newcommand{\bT}{\mathbf{T}}
\newcommand{\bG}{\mathbf{G}}
\newcommand{\bu}{\mathbf{u}}
\newcommand{\bs}{\mathbf{s}}
\newcommand{\br}{\mathbf{r}}
\newcommand{\bx}{\mathbf{x}}
\newcommand{\bR}{\mathbf{R}}
\newcommand{\bpsi}{\boldsymbol{\psi}}
\newcommand{\tp}{^{\mkern-1.5mu\mathsf{T}}}

\title{\bfseries The Voxel \texorpdfstring{$T$}{T}-Matrix\\[4pt]
\large Discretising an elastic scatterer into cubes: the matched propagator,\\
and the accuracy of the result against exact Mie}
\author{}
\date{19 September 2026}

\begin{document}
\maketitle

\begin{abstract}
\noindent
A scatterer of arbitrary shape can be represented as a lattice of small cubes,
each carrying an analytic single-site $T$-matrix, coupled through the
elastodynamic propagator and solved by Foldy--Lax.  This note sets out that
construction, identifies the condition the propagator must satisfy for the
scheme to be consistent, and measures the accuracy that results.

The condition is a pairing.  A single-site $T$-matrix closes on a particular
functional of the internal field --- the value at the cube centre, or the
average over the cube --- and the propagator it is used with must close on the
same one.  The collocation closure requires the propagator averaged over the
\emph{receiver cell only}: one power of the cell form factor, not the two of a
Galerkin pairing and not the none of a point evaluation.

That average is given in closed form.  It is regular at every non-zero lattice
separation, its far field is a scalar multiple of the point propagator per
mode, and near contact it needs a quadrature order that follows the separation.

Accuracy is measured against the exact elastic Mie solution for a sphere, which
is the one arbiter that is both exact and laterally varying.  Two controls are
needed for the comparison to measure the method rather than the test: the
observation distance must reach the far field, and the staircase volume of the
voxelised body must be divided out.  With those in place the matched pairing is
$2.9$--$5.8\times$ closer to Mie than a point propagator in the Rayleigh
regime.  Beyond it the pairing does not help: at $ka = 0.5$ and $ka = 1.5$ the
averaged propagator is slightly \emph{worse} than the point one, and what limits
the method is the staircase representation of a curved boundary.

\emph{These numbers supersede an earlier set measured against an arbiter that
carried three compensating sign errors; \S\ref{sec:results} records what
changed.}
\end{abstract}

\tableofcontents

%=====================================================================
\section{The construction}\label{sec:construction}

\subsection{Foldy--Lax on a lattice of cubes}

Partition the scatterer into $N$ cubes of side $d$ on a regular lattice, with
centres $\br_m$.  Each cube is represented by a local $T$-matrix $\bT_{\rm loc}$
acting on a nine-component state --- three displacements and six Voigt strains,
\begin{equation}
  \bpsi = \big(u_z,\,u_x,\,u_y,\;\varepsilon_{zz},\,\varepsilon_{xx},\,
               \varepsilon_{yy},\,2\varepsilon_{xy},\,2\varepsilon_{zy},\,
               2\varepsilon_{zx}\big)\tp,
\end{equation}
which is the smallest state closed under the coupling: the source a cube
radiates depends on both the displacement and the strain it experiences.

The exciting field at cube $m$ is the incident field plus the field radiated by
every other cube,
\begin{equation}\label{eq:foldy}
  \bpsi_m = \bpsi^{\rm inc}_m + \sum_{n\neq m} \bP(\br_m-\br_n)\,\bT_{\rm loc}\,\bpsi_n,
  \qquad\text{that is}\qquad
  \big(\bI - \bP\bT\big)\bpsi = \bpsi^{\rm inc},
\end{equation}
with $\bP(\bs)$ the $9\times9$ propagator between cubes at separation $\bs$.
The composite $T$-matrix of the whole body follows by summing the sources,
$\bT_{\rm comp} = \sum_m \bT_{\rm loc}\bpsi_m$.

The self term is excluded from the sum: a cube's action on itself is contained
in $\bT_{\rm loc}$, which is where the singular part of the Green's tensor and
its static (Eshelby) content live.

\subsection{The propagator}

$\bP$ is assembled from the elastodynamic Green's tensor and its derivatives,
in blocks
\begin{equation}\label{eq:Pblocks}
  \bP = \begin{pmatrix} \bG & \mathbf{C}\\[2pt] \mathbf{H} & \mathbf{S}\end{pmatrix},
\end{equation}
where $\bG$ maps force to displacement, $\mathbf{S}$ maps stress source to
strain, and $\mathbf{C},\mathbf{H}$ are the cross couplings obtained by Voigt
contraction of $\partial G$ and $\partial\partial G$.  Written through the
Kupradze representation, every block is built from the two scalar Helmholtz
kernels $g_P$ and $g_S$ and their Cartesian derivative tensors, so the $P$ and
$S$ content is separable --- which \S\ref{sec:average} requires.

%=====================================================================
\section{The pairing condition}\label{sec:pairing}

\subsection{Two closures}

Volume-averaging the Lippmann--Schwinger equation with a uniform internal
ansatz gives
\begin{equation}\label{eq:double}
  \langle\nabla\bu\rangle = \langle\nabla\bu^0\rangle
   + \left[\frac{1}{V}\iint \partial\partial G(\bx-\bx')\,d\bx\,d\bx'\right]\Delta c\;E,
\end{equation}
a \emph{double} average, while evaluating the same equation at the cube centre
gives
\begin{equation}\label{eq:single}
  \nabla\bu(0) = \nabla\bu^0(0)
   + \left[\int_V \partial\partial G(-\bx')\,d\bx'\right]\Delta c\;E,
\end{equation}
a \emph{single} average.  Neither is wrong: one closes on the volume average of
the internal field, the other on its centre value.  For an ellipsoid they
coincide, because Eshelby's uniformity theorem makes the internal field
uniform; for a cube they differ.

\subsection{The condition}

A $T$-matrix built on \eqref{eq:single} must be used with a propagator that
also closes on the centre value.  The propagator between two cubes therefore
carries \emph{one} power of the cell form factor --- the average over the
receiver cell, with the source treated as a point --- and not two.  Writing the
cell indicator's transform,
\begin{equation}\label{eq:formfactor}
  \text{single average}\;\;\longleftrightarrow\;\;
    \operatorname{sinc}(k_xh)\operatorname{sinc}(k_yh)\operatorname{sinc}(k_zh),
  \qquad h = d/2,
\end{equation}
against the squared factor that a Galerkin pairing produces.  A point
propagator carries neither and is not a consistent partner for either closure.

\paragraph{Per mode.}  The form factor must be applied to each mode separately.
The $P$ and $S$ poles carry different $k_z$ at the same lateral wavenumber, so
a single scalar factor applied to the assembled $9\times9$ is incorrect.

%=====================================================================
\section{The single cell average}\label{sec:average}

\subsection{Definition and regularity}

For cubes at separation $\bs$ on the lattice, the matched propagator is
\begin{equation}\label{eq:avgdef}
  \langle\bP\rangle(\bs) = \frac{1}{d^3}\int_{|\mathbf{u}|_\infty\le h}
     \bP(\bs-\mathbf{u})\,d\mathbf{u}, \qquad h = d/2 .
\end{equation}
With a point source and a cubic receiver cell, the closest approach of the cell
to the source is its near face, so
\begin{equation}\label{eq:regular}
  |\bs-\mathbf{u}| \;\ge\; d/2 \qquad\text{for every lattice separation }\bs\neq0,
\end{equation}
and the integrand in \eqref{eq:avgdef} is regular everywhere on the domain.
The double average has no such property: its integrand reaches $r=0$ whenever
the cells touch.

\subsection{The far field, in closed form}

The midpoint error of a cell average is
\begin{equation}\label{eq:tail}
  \langle g\rangle - g = \frac{d^2}{24}\nabla^2 g + O(d^4),
\end{equation}
and for a Helmholtz kernel $\nabla^2 g_c = -\kappa_c^2 g_c$ exactly away from
the origin.  Derivatives commute with $\nabla^2$, so the same holds for every
derivative tensor and hence for every block of \eqref{eq:Pblocks}.  Therefore,
away from contact,
\begin{equation}\label{eq:scalartail}
  \langle\bP\rangle = \Big(1 - \frac{\kappa_c^2 d^2}{24}\Big)\bP + O(d^4)
  \qquad\text{per mode } c\in\{P,S\},
\end{equation}
the two modes carrying their own $\kappa$ and therefore not combinable before
the average is taken.

\subsection{Quadrature near contact}

Inside the near shell \eqref{eq:scalartail} is not accurate enough and
\eqref{eq:avgdef} is integrated directly.  The integrand is regular but peaked
near the face of the cell closest to the source, and a Gauss rule converges
slowly against a nearby peak.  The difficulty is a property of the separation
alone,
\begin{equation}\label{eq:qparam}
  q(\bs) = \frac{\big|\max(|\bs|-h,\,0)\big|}{d},
\end{equation}
componentwise in the clamp, which takes the values $0.5$ for a face neighbour,
$1/\sqrt2$ for an edge neighbour and $\sqrt3/2$ for a corner one.  Convergence
of a product Gauss rule against a $24$-point reference:

\begin{center}
\begin{tabular}{@{}lrrrr@{}}
\toprule
neighbour & $g=4$ & $g=8$ & $g=12$ & $g=16$\\
\midrule
face $(q=0.50)$   & $1.8\times10^{-2}$ & $2.9\times10^{-5}$ & $3.6\times10^{-8}$ & $1.4\times10^{-10}$\\
edge $(q=0.71)$   & $6.5\times10^{-4}$ & $9.0\times10^{-8}$ & $1.6\times10^{-10}$ & ---\\
corner $(q=0.87)$ & $6.1\times10^{-5}$ & $2.3\times10^{-10}$ & --- & ---\\
\bottomrule
\end{tabular}
\captionof{table}{The order should follow $q$: $16$ points per axis for face
neighbours, $12$ for edge, $8$ for corner, and the base order elsewhere.  Only
six and twelve separations respectively are expensive, so the adaptive rule
costs almost nothing.}
\end{center}

\noindent
Two independent evaluations of \eqref{eq:avgdef} --- direct quadrature and the
closed form \eqref{eq:scalartail} --- agree to $O(d^4)$, halving $d$ reducing
their difference by sixteen ($5.3\times10^{-4}$, $3.4\times10^{-5}$,
$2.1\times10^{-6}$ at $d = 0.2,\,0.1,\,0.05$), and the average reduces to the
point propagator at second order as \eqref{eq:tail} requires.

\subsection{Cost}

$\bP$ depends only on the separation, so on a regular lattice there are
$O(n^3)$ distinct blocks rather than $O(N^2)$ pairs.  Building the kernel once
per distinct separation --- which is what an FFT-accelerated solver does
anyway --- makes the averaged propagator affordable: the few near-contact
separations needing high-order quadrature are evaluated once, whatever the cell
count.

%=====================================================================
\section{Measuring the accuracy}\label{sec:measure}

The exact elastic Mie solution for a homogeneous sphere is the natural arbiter:
it is exact, and unlike a layered reference it is laterally varying.  Two
things must be controlled for the comparison to measure the discretisation
rather than the test.

\subsection{The observation distance}

If the scattered field is compared through a far-field asymptotic on one side
and an exact expression on the other, the neglected $1/(kr)$ term must be
small compared with the quantity being measured.  It is not a weak condition.
Measured at $ka = 0.1$ with the sphere divided $4$ per edge:

\begin{center}
\begin{tabular}{@{}rrr@{}}
\toprule
$r/a$ & $k_S r$ & volume-corrected error\\
\midrule
$100$ & $10$ & $0.5043$\\
$500$ & $50$ & $0.0575$\\
$5\times10^{3}$ & $500$ & $0.0075$\\
$5\times10^{4}$ & $5000$ & $0.0068$\\
\bottomrule
\end{tabular}
\captionof{table}{The comparison is usable only from $k_Sr\sim500$; at
$k_Sr=50$ it overstates the error sixfold.  It is not fully converged even at
the top: the last decade still moves the answer by $10\%$, which sets the
precision of every small number in \S\ref{sec:results}.  The earlier version of
this table read $0.0088$ at both of the last two rows and so looked converged;
correcting the arbiter made the errors smaller without making the harness
better, and the residual $1/(k r)$ term is now a visible fraction of them.}
\end{center}

\subsection{The staircase}

A cubic lattice does not fill a sphere, and its coverage is not monotonic in
the number of cells per edge.  For $n_{\rm sub} = 3,4,6$ the voxel volume
exceeds the sphere's by $26\%$, $4.7\%$ and $17\%$, and the raw error tracks
that dilution to within a factor $1.4$.  Two consequences follow.

First, the raw error must be divided by the volume ratio before it says
anything about the propagator.  Second, a refinement ladder in $n_{\rm sub}$
measures the geometry, not the method: the error moves non-monotonically
because the dilution does.  The comparison that isolates the discretisation is
at \emph{fixed} geometry --- the same voxel set, the same staircase, one
propagator exchanged for the other --- and its stability across resolutions is
what makes it meaningful.

%=====================================================================
\section{Results}\label{sec:results}

\begin{quote}
\small\itshape
What changed, and why.  Every number in this section is a comparison between two
objects that were \emph{both} wrong when it was first measured.  The exact Mie
solution carried a spurious $(-1)^n$ in its multipole sum, inverting the
density/dipole channel; the Foldy--Lax far field carried a sign error in its
force term; and the contrast-extraction route carried a third error that
cancelled the first, which is how all three survived more than a thousand
tests.  They were found by an impedance-null test --- $\delta\rho/\rho = +
\varepsilon$ with $\delta M/M = -\varepsilon$ leaves $Z=\sqrt{\rho M}$ unchanged
and therefore cannot backscatter --- and fixed together; both solvers now agree
with an independent Born calculation to $0.9997$--$1.0000$ across $k_Pa =
0.18$ to $1.44$.

Three conclusions of the earlier version do not survive the correction, and are
stated here rather than quietly amended.  \textbf{(i)} The pairing improvement at
$ka = 0.5$ reversed sign: the averaged propagator was reported $1.15$--$1.19
\times$ better and is $0.71$--$0.84\times$, that is worse.  \textbf{(ii)} The
Rayleigh improvement was reported stable to $8.8\%$ and has a spread of $64\%$.
\textbf{(iii)} The resonance error roughly halved, from $52\%$ to $22$--$28\%$,
and refinement from $n_{\rm sub} = 6$ to $8$ now \emph{improves} it where it
previously made it worse.  What survives is the Rayleigh result and the whole of
\S\ref{sec:pairing}--\S\ref{sec:average}, which concern the propagator alone and
never passed through the arbiter.
\end{quote}

Volume-corrected far-field error against exact Mie, for a sphere with
$\Delta\lambda = 2$, $\Delta\mu = 1$, $\Delta\rho = 0.1$ on a background
$\alpha=5$, $\beta=3$, $\rho=2.5$:

\begin{center}
\begin{tabular}{@{}rrrrr@{}}
\toprule
$ka$ & $n_{\rm sub}$ & point propagator & matched average & ratio\\
\midrule
$0.1$ & $3$ & $0.0055$ & $\mathbf{0.0019}$ & $2.86$\\
      & $4$ & $0.0068$ & $\mathbf{0.0014}$ & $5.02$\\
      & $6$ & $0.0086$ & $\mathbf{0.0015}$ & $5.76$\\
\midrule
$0.5$ & $3$ & $\mathbf{0.0306}$ & $0.0365$ & $0.84$\\
      & $4$ & $\mathbf{0.0141}$ & $0.0181$ & $0.78$\\
      & $6$ & $\mathbf{0.0198}$ & $0.0278$ & $0.71$\\
\midrule
$1.5$ & $6$ & $\mathbf{0.2780}$ & $0.2851$ & $0.97$\\
      & $8$ & $\mathbf{0.2170}$ & $0.2247$ & $0.97$\\
\bottomrule
\end{tabular}
\captionof{table}{Bold is the better of the pair.  Matching the pairing helps in
the Rayleigh regime and nowhere else: at $ka = 0.5$ the averaged propagator is
$20$--$40\%$ \emph{worse}, consistently across resolutions, and at $ka = 1.5$ it
is $3\%$ worse.  The Rayleigh improvement is also not as stable as it looked ---
$2.86$ to $5.76$, a spread of $64\%$, driven by the coarsest geometry.}
\end{center}

\subsection{The three regimes}

\paragraph{Rayleigh, $ka\lesssim0.1$.}  The matched average reaches
$0.14$--$0.19\%$, at or below the shape-factor floor of representing a sphere by
cubes, against $0.55$--$0.86\%$ for the point propagator.  This is where the
pairing argument of \S\ref{sec:pairing} is confirmed, and it is the one regime
in which it is.  The improvement is not uniform: $2.86\times$ at
$n_{\rm sub}=3$ against $5.76\times$ at $n_{\rm sub}=6$.  The coarsest case is
also the most diluted ($26\%$), so the spread is plausibly geometric rather than
a property of the propagator --- but that is an explanation offered, not one
tested, and the earlier claim of stability across resolutions is withdrawn.

\paragraph{Transition, $ka\sim0.5$.}  \textbf{The pairing stops helping and
starts hurting}: $0.71$--$0.84\times$, consistently across all three
resolutions.  The effect is systematic, not scatter, and it is larger than the
harness uncertainty of the previous subsection.  Two readings are available and
this note does not choose between them.  Either the collocation closure that
\S\ref{sec:pairing} matches is itself losing validity as $ka_{\rm sub}$ grows,
so that matching it matches the wrong thing; or the cell average is correct and
is being scored against a staircase error that happens to cancel part of the
point propagator's own.  Distinguishing them needs a body whose boundary the
lattice resolves exactly --- a cube, not a sphere --- which is a different
measurement from this one.

\paragraph{Resonance, $ka\gtrsim1$.}  The error is $22$--$28\%$, and the
propagator is not what is wrong: the two propagators differ by $3\%$ where the
error is $25\%$.  The sub-cell $T$-matrices remain inside their Rayleigh
validity ($ka_{\rm sub} = 0.25$ at $n_{\rm sub}=6$), and refining from
$n_{\rm sub}=6$ to $8$ now improves the error from $27.8\%$ to $21.7\%$ as the
dilution falls from $17\%$ to $4.3\%$ --- so what the ladder is tracking here is
still the staircase, not the method.  What fails is the representation of a
curved boundary at a wavelength comparable with the radius.  That is a limit of
the discretisation itself, not of the propagator within it.

%=====================================================================
\section{Summary}

\begin{enumerate}
\item Partition the body into cubes; give each the analytic single-site
      $T$-matrix for its contrast.
\item Couple them with the propagator averaged over the \emph{receiver} cell,
      \eqref{eq:avgdef}, applied per mode --- matching the collocation closure
      of the single-site $T$-matrix.
\item Use the closed form \eqref{eq:scalartail} beyond a near shell and direct
      quadrature inside it, with the order chosen from \eqref{eq:qparam}.
\item Solve \eqref{eq:foldy}; build the kernel once per distinct separation.
\item When comparing against an exact reference, reach the far field and
      divide out the staircase volume; compare at fixed geometry.
\end{enumerate}

\noindent
The method is accurate to a fraction of a percent in the Rayleigh regime,
degrades to one to four percent through the transition, and is geometry limited
beyond $ka\sim1$ at some twenty-five percent, where a representation that
follows the boundary rather than staircasing it is required instead.

\noindent
The pairing of step~2 is what the Rayleigh figure rests on, and it is
established only there.  At $ka = 0.5$ and above the averaged propagator is
measurably worse than the point one, and until that is understood step~2 should
be read as a Rayleigh-regime recommendation rather than a general one.

\end{document}
